\documentclass[10pt,a4paper]{amsart}
\usepackage[margin=19mm]{geometry}
\usepackage{amsmath,amssymb,mathtools}
\usepackage[scr=rsfs,cal=euler]{mathalfa}
\usepackage{xfrac}
\usepackage[unicode,hidelinks,bookmarks=false]{hyperref}
\newcommand{\ie}{i.e.\ }
\newcommand{\cf}{cf.\ }
\newcommand{\resp}{respectively\ }
\newcommand{\Subsection}{\subsection}
\providecommand{\nobreakdash}{\nobreak\hbox{-}}
\setlength{\emergencystretch}{2em}
\multlinegap0pt
\usepackage{mathtools}
\usepackage{amssymb}
\usepackage{tcolorbox}
\usepackage{tikz}
\usepackage{float}
\newtheorem{lemma}{Lemma}
\title{Local Limits of Small World Networks}
\author{Yeganeh Alimohammadi, Senem I\c{s}\i k, Amin Saberi}
\date{Source: arXiv:2501.11226v2 (CC BY 4.0)}
\begin{document}
\maketitle

\noindent\emph{Source and licence.} Local Limits of Small World Networks. Version arXiv:2501.11226v2 (CC BY 4.0), distributed by arXiv under the Creative Commons Attribution 4.0 International licence, http://creativecommons.org/licenses/by/4.0/, as recorded in the arXiv licence metadata for this exact version and shown on the abstract page https://arxiv.org/abs/2501.11226v2. Full licence notice: \texttt{licenses/arxiv-2501.11226-CC-BY-4.0.txt}.\par\smallskip

\noindent\textit{Editorial scope.} Counted unit: the article's lemma lem:critical-poisson-input (source0.tex lines 2162--2175). The article's complete original proof of it is delivered with it (source0.tex lines 3107--3224, one \texttt{proof} environment of 118 lines, which the article places away from the statement). No mathematics is written, repaired or re-derived: the statement and the proof are the article's own lines, in the article's own notation, and no retained line is substituted or re-flowed.  Retained uncounted as the source-local context the proof needs: lines 2100--2114.  Nothing was omitted from any retained span, so the package carries no omission record.  No retained line contains a citation, so the wrapper carries no bibliography and no bibliography entry.  No retained line contains a figure environment, an image-inclusion command or drawing code, so no image is delivered and the package declares no asset dependency.  Editorial formatting only, in addition: an AMS \texttt{amsart} preamble that restates the article's own declarations and macros used by the retained lines, namely the environments \texttt{lemma} on separate counters, together with the standard packages those lines need.  The retained environments print in this document as Lemma 1 in order of appearance, whereas the article prints the same items with the numbering of its own full document.  The complete native source of the article as distributed in the arXiv e-print for this exact version is bundled unchanged as \texttt{sources/171822/source0.tex}.
\begin{lemma}[Uniform normalization at criticality]\label{lem:uniform-normalization}
For every \(u\in V(G_n)\),
\[
D_n(u)\ge \log n+O(1).
\]
Moreover, if \(u\in V(G_n)\) has lattice distance at least
\[
d_n:=\left\lfloor \frac{n}{\log n}\right\rfloor
\]
from the boundary of the \(n\times n\) grid, then
\[
D_n(u)=4\log n+O(\log\log n),
\]
uniformly in \(u\).
\end{lemma}

\begin{lemma}[Critical limit distribution for incoming shortcuts is Poisson]\label{lem:critical-poisson-input}
Fix \(i\ge 1\). Let the nodes revealed before the \(i\)-th step of the BFS exploration started from some root, \(S_i\), be a fixed finite set of nodes,
with \(|S_i|=O(1)\), and let \(u\notin S_i\) satisfy
\(
S_i\cup\{u\}\subseteq G_{n-d_n}.
\) Define
\(
\mathrm{In}(u,V(G_n)\setminus S_i)
\)
to be the number of incoming shortcuts to \(u\) from nodes from \(V(G_n)\setminus S_i\). Then, 
\[
\mathrm{In}(u,V(G_n)\setminus S_i)\xrightarrow{d}\mathrm{Poi}(q).
\]
\end{lemma}

\begin{proof}[Proof of Lemma~\ref{lem:critical-poisson-input}]
For each \(v\in V(G_n)\), let \(\mathrm{out}_{v,i}\in\{0,\dots,q\}\) denote the
number of outgoing shortcuts from \(v\) that have already been revealed before
step \(i\). For each \(v\notin S_i\cup\{u\}\), let \(Y_{v,i}\) denote the number
of still-unrevealed outgoing shortcuts from \(v\) that land at \(u\). Then, we have
\[
Y_{v,i}\sim \mathrm{Bin}\!\left(q-\mathrm{out}_{v,i},\,p_{v,i}\right),
\qquad
p_{v,i}:=\frac{d(v,u)^{-2}}{D_n(v)},
\]
and the family \(\{Y_{v,i}\}_{v\notin S_i\cup\{u\}}\) is independent. Hence, the number of incoming shortcuts satsify
\[
\mathrm{In}(u,V(G_n)\setminus S_i)
=\sum_{v\notin S_i\cup\{u\}} Y_{v,i}.
\]
Define
\[
\lambda_n:=\sum_{v\notin S_i\cup\{u\}} (q-\mathrm{out}_{v,i})p_{v,i},
\]
giving the mean of \(\mathrm{In}(u,V(G_n)\setminus S_i)\). We
first show that \(\lambda_n=q+o(1)\).
Set
\[
b_n:=\Big\lfloor \frac{d_n}{2}\Big\rfloor \text{ for } 
d_n:=\left\lfloor \frac{n}{\log n}\right\rfloor,
\qquad
N_n:=\{w\in V(G_n): d(w,u)\le b_n\}.
\]
Since \(u\) is at lattice distance at least \(d_n\) from the boundary, every
\(w\in N_n\) is at lattice distance at least \(d_n-b_n\ge b_n\) from the
boundary. Repeating the proof of Lemma~\ref{lem:uniform-normalization} with
\(b_n\) in place of \(d_n\), we obtain
\[
D_n(w)=4\log n+O(\log\log n)
\qquad\text{uniformly for }w\in N_n.
\]
Therefore, we get
\[
\sum_{w\in N_n} q\,p_{w,i}
=\frac{q}{4\log n+O(\log\log n)}\sum_{j=1}^{b_n}4j\cdot j^{-2}
=q\,\frac{4\log b_n+O(1)}{4\log n+O(\log\log n)}
=q+o(1),
\]
since \(\log b_n/\log n\to 1\).

For the complement \(N_n^c\), Lemma~\ref{lem:uniform-normalization} gives the
uniform lower bound
\[
D_n(w)\ge \log n+O(1)
\qquad\text{for all }w\in V(G_n),
\]
giving
\[
\sum_{w\notin N_n} q\,p_{w,i}
\le \frac{q}{\log n+O(1)}\sum_{j>b_n}4j\cdot j^{-2}
=O\!\left(\frac{\log(n/b_n)}{\log n}\right)
=o(1).
\]

It remains to control the correction coming from the factors
\(\mathrm{out}_{v,i}\). Because $|S_i| = O(1)$ by assumption, we have
\(
\sum_{v\in V(G_n)} \mathrm{out}_{v,i} = C,
\)
for some constant $C$.
Using again \(D_n(v)\ge \log n+O(1)\) and \(d(v,u)\ge 1\), we get
\[
p_{v,i}\le \frac{1}{\log n+O(1)},
\qquad
\sum_{v\notin S_i\cup\{u\}} \mathrm{out}_{v,i}p_{v,i}
\le \frac{C}{\log n+O(1)}
=o(1).
\]
Combining the near contribution, the far contribution, and this correction term
yields
\(
\lambda_n=q+o(1).
\)
We now identify the full conditional law. Here and below,
\(\mathcal L(X)\) denotes the law of \(X\). Since
\(\mathrm{In}(u,V(G_n)\setminus S_i)\) is a sum of independent Bernoulli random
variables, Le Cam's inequality gives
\[
d_{\mathrm{TV}}\!\left(
\mathcal L\bigl(\mathrm{In}(u,V(G_n)\setminus S_i)\bigr),
\,\mathrm{Poi}(\lambda_n)\right)
\le
2\sum_{v\notin S_i\cup\{u\}} (q-\mathrm{out}_{v,i})p_{v,i}^2.
\]
Using again \(D_n(v)\ge \log n+O(1)\), we obtain
\(
p_{v,i}^2\le \frac{d(v,u)^{-4}}{(\log n+O(1))^2}.
\)
Therefore, we have
\[
\sum_{v\notin S_i\cup\{u\}} (q-\mathrm{out}_{v,i})p_{v,i}^2
\le
\frac{q}{(\log n+O(1))^2}\sum_{j\ge 1}4j\cdot j^{-4}
=O\!\left(\frac{1}{\log^2 n}\right)\to 0,
\]
and
\[
d_{\mathrm{TV}}\!\left(
\mathcal L\bigl(\mathrm{In}(u,V(G_n)\setminus S_i)\bigr).
\,\mathrm{Poi}(\lambda_n)\right)\to 0
\]
Since \(\lambda_n\to q\), we
also have
\(
d_{\mathrm{TV}}\bigl(\mathrm{Poi}(\lambda_n),\mathrm{Poi}(q)\bigr)\to 0.
\)
By the triangle inequality,
\[
d_{\mathrm{TV}}\!\left(
\mathcal L\bigl(\mathrm{In}(u,V(G_n)\setminus S_i) \bigr),
\,\mathrm{Poi}(q)\right)\to 0.
\]
\end{proof}

\end{document}
